\begin{abstract}
Although large-scale unsupervised language models (LMs) acquire extensive world knowledge and display some reasoning capabilities, exercising precise command over their behavior remains challenging due to the absence of supervision during training. Common techniques to introduce such controllability involve gathering human judgments about the comparative quality of generated outputs and subsequently fine-tuning the base LM to better reflect these preferences, frequently utilizing reinforcement learning from human feedback (RLHF). RLHF, however, is typically cumbersome and prone to instability, as it necessitates first constructing a reward model representative of human preferences and then applying reinforcement learning to further adjust the LM, optimizing for the surrogate reward without deviating excessively from its original behavior. In this work, we propose a novel reparameterization of the reward model used in RLHF, which permits direct computation of the associated optimal policy via a closed-form solution. This approach allows us to reformulate the canonical RLHF objective so that it can be addressed using a straightforward classification loss. Our resulting framework, termed \textit{Direct Preference Optimization} (DPO), is robust, efficient, and requires minimal computational resources, obviating the need for LM sampling during the fine-tuning phase and reducing reliance on extensive hyperparameter search. Through empirical evaluation, we demonstrate that DPO fine-tunes LMs to reflect human preferences with performance comparable to or surpassing traditional approaches. Significantly, DPO not only provides better sentiment control than PPO-based RLHF but also equals or outperforms it in response quality across tasks such as summarization and single-turn dialogue, all while greatly simplifying both implementation and training.
\end{abstract}

\section{Introduction}
Large unsupervised language models (LMs) trained on massive corpora have demonstrated remarkable and often unexpected competencies~\citep{chowdhery2022palm, brown2020language, touvron2023llama,bubeck2023sparks}. Nevertheless, these systems learn from human-produced data encompassing a broad spectrum of intentions, values, and expertise. As a result, not every behavior or skill reflected in the training data is appropriate to emulate; for instance, an AI coding assistant should ideally \textit{recognize} common programming mistakes to fix them, yet, during code synthesis, the assistant ought to prefer the (sometimes uncommon) instances of expert coding present in its training data. Similarly, a language model might need to be \textit{informed} about widespread misconceptions—such as one held by half the population—but we do not want the model to propagate this falsehood in half of its relevant outputs. In essence, the ability to \emph{select and control} the model’s outputs—choosing desired behaviors from its extensive repertoire of \textit{knowledge and skills}—is fundamental to achieving safety, reliability, and controllability in AI systems \citep{ouyang2022training}. Prevailing approaches tend to align LMs with human preferences using reinforcement learning (RL). In this work, however, we reveal that the RL-based objective often targeted in these methods can be solved exactly through a straightforward binary cross-entropy loss, which offers significant simplifications to the process of learning from preferences.

Generally, present-day techniques endow LMs with targeted behaviors via curated datasets of human preferences that define desirable and safe model conduct. This phase of preference alignment takes place subsequent to extensive unsupervised pre-training on large-scale textual datasets. Although the most direct technique for preference learning involves supervised fine-tuning on human-authored examples of optimal responses, the current leading methods utilize reinforcement learning from human (or AI) feedback (RLHF/RLAIF; \citep{christiano2017deep,bai2022constitutional}). In RLHF, a reward model is trained using human preference data, and this reward model guides the policy optimization of the language model through RL so as to increase the frequency of desirable outputs, while keeping the learned policy close to its original distribution. Although RLHF yields models with exceptional performance in dialogue and coding, the procedure is substantially more elaborate than supervised learning: it typically requires training multiple models and performing policy sampling within the RL training loop, which demands considerable computational resources.

This work demonstrates a method for optimizing a language model in accordance with human preferences, dispensing with the need for explicit reward modeling or reinforcement learning techniques. We introduce \textit{{Direct Preference Optimization} (DPO)}, an approach that implicitly targets the same objective as standard RLHF algorithms—namely, reward maximization subject to a KL-divergence regularization—but is easier to implement and train. At a high level, DPO updates the model by raising the relative log likelihood of responses that are preferred over those that are not, while employing a per-example, adaptive importance weight. This mechanism addresses the degeneration issues observed with a simple probability ratio objective. In common with prior methods, DPO assumes a theoretical preference model, such as the Bradley-Terry model (\cite{bradley1952rankanalysis}), which quantifies the consistency between a reward function and observed preference judgments. In contrast to prior approaches, which use the preference model to construct a loss function for training a reward model followed by optimization of the policy according to this reward, DPO performs a change of variables to define the preference loss directly in terms of the policy. Utilizing a collection of human preference comparisons among model responses, DPO trains a policy using a binary cross entropy loss, resulting in an optimal policy under an implicit reward derived from the preference data.

Our principal contribution is Direct Preference Optimization (DPO), an RL-free, straightforward method for language model training from preferences. Experimental evaluations demonstrate that DPO achieves comparable or superior performance to established techniques, including RLHF methods based on PPO, across applications such as sentiment control, summarization, and dialogue with models as large as 6B parameters.

\section{Related Work}

Self-supervised language models that are scaled up can perform certain tasks with zero-shot learning \citep{radford2019language} or by employing few-shot prompt-based techniques \citep{gpt3,megatron,chowdhery2022palm}. Nevertheless, their efficacy on specific applications and adherence to user intentions are markedly enhanced when they are further trained on datasets consisting of explicit instructions paired with human-authored responses \citep{mishra-etal-2022-cross,sanh2022multitask,chung2022scaling,thoppilan2022lamda}. Through this process, termed `instruction-tuning', large language models (LLMs) gain the ability to follow novel instructions not present during tuning, thereby increasing their overall utility \citep{chung2022scaling}. Despite its achievements, it is often more feasible to gather \textit{relative} assessments of response quality from non-experts than to obtain expert-constructed demonstrations. Consequently, subsequent research has employed datasets of human preferences to refine LLMs, thereby achieving improvements in translation \citep{kreutzer-etal-2018-reliability}, summarization \citep{stiennon2022learning,ziegler2020finetuning}, narrative generation \citep{ziegler2020finetuning}, and the ability to follow instructions \citep{ouyang2022training,ramamurthy2023is}. The standard approach consists of initially fitting a neural reward model to align with the recorded human preferences, often modeled using frameworks like the Bradley-Terry model \citep{bradley1952rankanalysis}. Next, the language model is adapted to maximize this reward via reinforcement learning algorithms, typically REINFORCE \citep{williams1992reinforce}, proximal policy optimization (PPO; \cite{schulman2017proximal}), or related methodologies \citep{ramamurthy2023is}. Parallel lines of research use instruction-following LLMs, enhanced through human feedback, to synthesize further preference data tailored to specific aspects such as safety or non-harmfulness \citep{bai2022constitutional}, requiring only minimal human oversight via descriptive rubrics for annotation. This area blends two major traditions: reinforcement learning-based approaches to training language models for diverse objectives~\citep{Ranzato2015SequenceLT,paulus2018a,wu2018learning}, and more general strategies for learning from human feedback \citep{christiano2017deep,kupcsik2018learning}. However, despite the practical advantages of leveraging comparative human preferences, refining large language models through reinforcement learning continues to present considerable technical difficulties; in response, this work offers a theoretically sound method for directly optimizing relative preferences, thereby obviating the need for RL.

Beyond language applications, learning policies from preferences has been explored in both bandit and reinforcement learning paradigms, giving rise to multiple methodologies. In particular, contextual bandit approaches that utilize preferences or ordinal rankings of actions in lieu of reward signals are referred to as contextual dueling bandits (CDB; \cite{yue2012karmed,dudik2015contextual}). When absolute reward feedback is unavailable, theoretical treatments of CDBs replace the conventional optimal policy notion with the concept of a \textit{von Neumann winner}—a policy that maintains an expected winning rate of at least 50\% against every alternative policy \citep{dudik2015contextual}. It should be noted that in the CDB framework, preference feedback is provided sequentially (online), in contrast to preference-based policy learning from humans, which typically utilizes a fixed, offline collection of action pairs annotated with preference labels \citep{yan2022human}. Similarly, \textit{preference-based RL} (PbRL) investigates the use of binary preference signals generated by an unobserved 'scoring' function as substitutes for explicit rewards \citep{BusaFekete2014,ruiz2023dueling}. A variety of algorithms address PbRL, some of which permit reusing off-policy preference data; most commonly, they involve first estimating the hidden scoring function—i.e., constructing a reward model—followed by an explicit optimization of that model \citep{jain2013learning,BusaFekete2014,christiano2017deep,sadigh2017active,kupcsik2018learning}. In contrast, our work introduces a single-stage approach in which the policy is optimized directly to align with the observed preferences.

\section{Preliminaries}\label{section:prelims}

Here, we summarize the RLHF methodology outlined by \citeauthor{ziegler2020finetuning} and further discussed in subsequent work \citep{stiennon2022learning, bai2022training, ouyang2022training}. The typical RLHF process is comprised of three main components: 1) supervised fine-tuning (SFT); 2) preference data collection and reward modeling; and 3) reinforcement learning optimization.

\textbf{SFT}: The RLHF pipeline commonly begins by adapting a pre-trained language model via supervised learning, leveraging curated data relevant to the specific downstream application (e.g., dialogue, summarization). This produces a model referred to as $\pi^\text{SFT}$.

\textbf{Reward Modeling Phase}: During this phase, the SFT model is given prompts $x$ and generates candidate completions, yielding answer pairs $(y_1, y_2)\sim \pi^\text{SFT}(y \mid x)$. Human annotators are then presented with these pairs and indicate their preference, denoted as $y_w\succ y_l \mid x$, where $y_w$ stands for the favored response and $y_l$ for the less favored one. The underlying assumption is that these preferences are sampled according to an unknown reward function $r^*(y, x)$, which remains inaccessible. Preferences can be formalized using several probabilistic models, with the Bradley-Terry (BT) \cite{bradley1952rankanalysis} model being a widely adopted example (though more elaborate schemes such as Plackett-Luce models \citep{plackett1975analysis, luce2012individual} are also compatible, particularly if ranked orderings over multiple responses are collected). The BT model formalizes the probability distribution over preferences $p^*$ as:

\begin{equation}\label{eq:bradley-terry}
    p^*(y_1\succ y_2 \mid x)=\frac{\exp\left(r^*(x, y_1)\right)}{\exp\left(r^*(x, y_1)\right) + \exp\left(r^*(x, y_2)\right)}.
\end{equation}

Given a fixed dataset of comparisons $\mathcal{D}=\bigl\{x^{(i)}, y_w^{(i)}, y_l^{(i)}\bigr\}_{i=1}^N$ sampled according to $p^*$, a reward function $r_{\phi}(x, y)$ can be parametrized, and its parameters estimated by maximizing the likelihood. By viewing this setup as a binary classification problem, we can write the negative log-likelihood loss as follows:

\begin{equation}\label{eq:reward_model}
    \mathcal{L}_R(r_{\phi}, \mathcal{D}) = -\mathbb{E}_{(x, y_w, y_l)\sim \mathcal{D}}\bigl[\log \sigma(r_{\phi}(x, y_w)- r_{\phi}(x, y_l))\bigr]
\end{equation}

where $\sigma$ denotes the logistic sigmoid. When applied to language models, $r_{\phi}(x, y)$ is frequently initialized from the SFT model $\pi^\text{SFT}(y \mid x)$, augmented by an additional linear head on top of the last transformer layer, yielding a single scalar output representing the reward \cite{ziegler2020finetuning}. Previous studies often normalize reward values to reduce variance, imposing $\mathbb{E}_{x,y\sim \mathcal{D}}\left[r_\phi(x, y)\right] = 0$ over all $x$.

\textbf{RL Fine-Tuning Phase}: In the subsequent reinforcement learning stage, the estimated reward model serves as the feedback mechanism for updating the language model. As in earlier research~\citep{jaques2017sequence, jaques2020human}, the problem is cast as

\begin{equation}\label{eq:RL}
\max_{\pi_{\theta}}  \mathbb{E}_{x\sim \mathcal{D}, y\sim \pi_{\theta}(y \mid x)}\bigl[r_{\phi}(x, y)\bigr] - \beta\mathbb{D}_{\textrm{KL}}\bigl[\pi_{\theta}(y\mid x)\mid \mid \pi_\text{ref}(y\mid x)\bigr],
\end{equation}

where $\beta$ determines how far the optimized policy can diverge from the reference policy $\pi_\text{ref}$, usually the initial SFT model $\pi^\text{SFT}$. In practical terms, $\pi_\theta$ is also commonly initialized with $\pi^\text{SFT}$. The KL-regularization is essential, as it constrains the policy to remain close to the distribution where the reward function is reliable and preserves diversity in generated responses, thus avoiding collapse to a single high-reward answer. Owing to the discreteness of language outputs, this objective is non-differentiable and generally optimized using reinforcement learning techniques. The established practice \citep{ziegler2020finetuning, stiennon2022learning, bai2022training, ouyang2022training} involves constructing the reward as ${r(x, y) = r_{\phi}(x, y) -\beta (\log \pi_{\theta}(y\mid x) - \log \pi_\text{ref}(y\mid x))}$, and applying PPO \cite{schulman2017proximal} for maximization.

\section{Direct Preference Optimization}\label{sec:DPO}

Recognizing the practical complexities in applying RL methods at the scale of modern language models, we seek a streamlined policy optimization method that directly utilizes preference data. Distinct from traditional RLHF pipelines, which require explicit reward modeling and RL-based maximization, our framework is built on a specific reward model structure allowing its optimal policy to be computed in closed form, thereby circumventing the RL optimization loop entirely. As elaborated below, our central observation is that an explicit, analytical relationship exists between reward models and their corresponding optimal policies, enabling us to translate a loss defined on rewards into a loss over policies. This transformation makes it unnecessary to learn a standalone reward model, while still capturing human preferences as described by models like the Bradley-Terry model. Ultimately, the policy network jointly embodies both the language generation mechanism and the (implicit) reward representation.

\textbf{Derivation of the DPO Objective.} We consider the standard reinforcement learning objective with a generic reward function $r$, as described in Eq.~\ref{eq:RL}, paralleling the approach taken in earlier work~\citep{peters2007reinforcement, peng2019advantage, korbak2022reinforcement, go2023aligning}. It is well-established that the optimal policy that maximizes expected reward under a KL constraint, as per Eq.~\ref{eq:RL}, takes the following form:

\begin{equation}\label{eq:op_policy}
    \pi_r(y\mid x) = \frac{1}{Z(x)}\pi_\text{ref}(y\mid x)\exp\left(\frac{1}{\beta}r(x, y)\right),
\end{equation}

where the normalization constant $Z(x) =\sum_{y}\pi_\text{ref}(y\mid x)\exp\left(\frac{1}{\beta}r(x, y)\right)$ acts as a partition function. For a full step-by-step derivation, refer to Appendix \ref{app:derivation1}. Although a maximum likelihood estimator $r_\phi$ may be available as an approximation to the true reward $r^*$, directly computing $Z(x)$ is typically computationally prohibitive~\citep{korbak2022reinforcement, go2023aligning}, limiting the practical usability of this form. Nonetheless, it is possible to rewrite Eq.~\ref{eq:op_policy} such that the reward function is expressed in terms of the optimal policy $\pi_r$, the reference distribution $\pi_\text{ref}$, and the unknown normalization $Z(\cdot)$. By applying the logarithm to both sides of Eq.~\ref{eq:op_policy} and re-arranging terms, we arrive at:

\begin{equation}\label{eq:main_eq}
    r(x,y) =\beta \log \frac{\pi_r(y\mid x)}{\pi_\text{ref}(y\mid x)} + \beta \log Z(x).
\end{equation}

Applying this alternative formulation to the true reward $r^*$ and its associated optimal policy $\pi^*$, we note that the Bradley-Terry model requires only the difference in rewards for a pair of outputs: ${p^*(y_1 \succ y_2 \mid x) = \sigma(r^*(x, y_1) - r^*(x, y_2))}$. Plugging the form in Eq.~\ref{eq:main_eq} for $r^*(x,y)$ into the preference relation in Eq.~\ref{eq:bradley-terry}, the partition function terms cancel, resulting in a human preference model expressed solely in terms of the optimal policy $\pi^*$ and the reference policy $\pi_\text{ref}$. Thus, under the Bradley-Terry setting, the optimal RLHF policy $\pi^*$ adheres to the preference model:

\begin{equation}\label{eq:objective}
    p^*(y_1\succ y_2 \mid x)=\frac{1}{1 + \exp\left(\beta \log \frac{\pi^*(y_2\mid x)}{\pi_\text{ref}(y_2\mid x)} - \beta \log \frac{\pi^*(y_1\mid x)}{\pi_\text{ref}(y_1\mid x)}\right)}
\end{equation}

Details of this derivation are given in Appendix~\ref{app:derivation2}. Although Eq.~\ref{eq:objective} is instantiated for the Bradley-Terry model, analogous formulations can be established for the Plackett-Luce family~\citep{plackett1975analysis, luce2012individual}; see Appendix~\ref{app:plackett_luce_models} for those derivations.

Having represented the probability of observed human preferences through the optimal policy rather than the reward model, we can now devise a maximum likelihood estimation problem for a parametric policy $\pi_\theta$. By analogy to the objective used in reward modeling (cf. Eq.~\ref{eq:reward_model}), our new policy-learning objective becomes:

\begin{equation}\label{eq:optimum_model}
    \mathcal{L}_\text{DPO}(\pi_{\theta}; \pi_\text{ref}) = -\mathbb{E}_{(x, y_w, y_l)\sim \mathcal{D}}\left[\log \sigma \left(\beta \log \frac{\pi_{\theta}(y_w\mid x)}{\pi_\text{ref}(

In this framework, we infer an implicit reward through an alternative form of parameterization, where the optimal policy coincides with $\pi_\theta$. Furthermore, because our methodology amounts to learning a reparameterized Bradley-Terry model, it inherits advantageous theoretical attributes, including consistency, given appropriate conditions on the distribution of preference data \cite{bong2022generalized}. Additional theoretical considerations regarding DPO and its relation to similar methods are explored in Section~\ref{sec:theory}.

\textbf{What is the effect of the DPO update?} To gain a deeper, mechanistic perspective on DPO, we examine the gradient of the objective $\mathcal{L}_\text{DPO}$. The derivative with respect to $\theta$ can be expressed as:

\begin{multline*}\label{eq:gradient}
    \nabla_\theta \mathcal{L}_\text{DPO}(\pi_\theta;\pi_\text{ref}) = \\ -\beta\mathbb{E}_{(x, y_w, y_l) \sim \mathcal{D}} \bigg[\underbrace{\sigma(\hat{r}_\theta(x, y_l) - \hat{r}_\theta (x, y_w))}_\text{higher weight when reward estimate is wrong}\bigg[\underbrace{\nabla_\theta\log \pi(y_w \mid x)}_\text{increase likelihood of $y_w$} - \underbrace{\nabla_\theta\log\pi(y_l \mid x)}_\text{decrease likelihood of $y_l$}\bigg]\bigg],
\end{multline*}

Here, the reward function $\hat{r}_\theta(x, y) = \beta \log \frac{\pi_\theta(y \mid x)}{\pi_\text{ref}(y \mid x)}$ is implicitly specified by the parameterized language model $\pi_\theta$ and the fixed reference $\pi_\text{ref}$ (see further discussion in Section~\ref{sec:theory}). Conceptually, taking the gradient of $\mathcal{L}_\text{DPO}$ acts to boost the probability assigned to the preferred outputs $y_w$ while reducing that of the less favored $y_l$. Crucially, the influence of each training example is modulated by the extent to which the implicit reward model $\hat{r}_\theta$ erroneously scores the dispreferred completion above the preferred one, with this error weighted by $\beta$—thus reflecting the misordering by the implicit reward model, subject to the intensity of the KL regularization. Empirically, this reweighting proves critical: omitting it, as in a simplistic version of the algorithm, often leads to model collapse, as detailed in Appendix Table~\ref{tab:unlikelihood_generations}.

\textbf{DPO overview.} The standard DPO procedure proceeds in the following manner: First, for each prompt $x$, two completions $y_1, y_2$ are sampled from the reference model $\pi_\text{ref}(\cdot \mid x)$, and are assigned preference labels by humans, thus forming the offline preference dataset $\mathcal{D} = \{x^{(i)}, y_w^{(i)}, y_l^{(i)}\}_{i=1}^N$. Second, the language model $\pi_\theta$ is trained by minimizing the DPO objective $\mathcal{L}_\text{DPO}$, making use of the fixed $\pi_\text{ref}$, the constructed dataset $\mathcal{D}$, and the desired parameter $\beta$.

In applied settings, researchers often rely on existing preference datasets rather than creating new completion samples and acquiring new human feedback. Since such datasets are generally produced using $\pi^\text{SFT}$, the reference policy is set to $\pi_\text{ref} = \pi^\text{SFT}$ whenever this model is accessible. In the absence of $\pi^\text{SFT}$, $\pi_\text{ref}$ is chosen by maximizing the likelihood over the set of preferred completions ${(x, y_w)}$ according to ${\pi_\text{ref} = \argmax_{\pi}\mathbb{E}_{x, y_w \sim \mathcal{D}}\left[\log \pi(y_w \mid x)\right]}$. This initialization procedure is designed to reduce the mismatch between the ideal but unavailable reference distribution and the actual $\pi_\text{ref}$ applied in DPO training. Appendix~\ref{app:implementation} presents additional details regarding the setup and choice of hyperparameters.

\section{Theoretical Analysis of DPO}
In what follows, we deepen the theoretical understanding of DPO, elucidate its conceptual underpinnings, and articulate its advantages in comparison to actor-critic techniques traditionally utilized for RLHF (such as PPO~\cite{schulman2017proximal}).

\label{sec:theory}

\subsection{Your Language Model Is Secretly a Reward Model}
DPO uniquely sidesteps the separate processes of reward model training and RL policy optimization by employing a unified maximum likelihood learning objective. Notably, the optimization in Eq.~\ref{eq:main_eq} can be viewed as a Bradley-Terry model where the reward is parameterized as $r^*(x, y) = \beta \log\frac{\pi^*_\theta(y \mid x)}{\pi_\text{ref}(y \mid x)}$. The model $\pi_{\theta}$ is then fit analogously to the standard reward modeling procedure (cf.~Eq.~\ref{eq:reward_model}) under a suitable change of variables. In this section, we formally establish the theoretical justification for this form of reparameterization, demonstrate its expressivity in representing arbitrary reward models, and prove its capacity to recover the optimal policy exactly. We begin by introducing a formal equivalence relation for reward functions.

\begin{definition}
We define two reward functions $r(x, y)$ and $r'(x, y)$ as equivalent if and only if ${r(x, y)-r'(x, y) = f(x)}$ for some function $f$.
\end{definition}

It is straightforward to verify that this notion of equivalence satisfies the properties of an equivalence relation, thereby partitioning the space of reward functions into distinct equivalence classes. This leads to the following two lemmas:

\begin{lemma}\label{lemma:same_prefrence}
Within the context of the Plackett-Luce model, and specifically the Bradley-Terry model, any two reward functions belonging to the same equivalence class induce identical preference distributions.
\end{lemma}

\begin{lemma}\label{lemma:same_policy}
    Within the constrained RL setting, any pair of reward functions from the same equivalence class yields the same optimal policy.
\end{lemma}

The arguments for these results are elementary and can be found in Appendix \ref{app:lemma1}. The first lemma points out a classic case of model under-specification common to the Plackett-Luce family \cite{plackett1975analysis}. This indeterminacy necessitates the introduction of extra identifiability conditions in order to make meaningful claims about the maximum likelihood estimates given by Eq. \ref{eq:reward_model} \cite{bong2022generalized}. The second lemma establishes that the optimal policy is invariant within a given equivalence class of reward functions; hence, for optimization purposes, it is sufficient to recover any member from the target equivalence class. The following theorem, whose proof is presented in Appendix~\ref{app:thm1}, formalizes this result:

\begin{theorem}\label{thm:main}
    Given mild technical conditions, each reward equivalence class associated with the Plackett-Luce (and, in particular, Bradley-Terry) framework can be characterized via the parameterization ${r(x, y) = \beta \log \frac{\pi(y\mid x)}{\pi_\text{ref}(y\mid x)}}$ for some model $\pi(y\mid x)$ and a fixed reference distribution $\pi_\text{ref}(y \mid x)$.
\end{theorem}

\begin{sproof}
    Take any reward function $r(x, y)$, which defines an optimal model $\pi_r(y \mid x)$ as per Eq. \ref{eq:op_policy}. Our goal is to demonstrate that there exists a reward function within the equivalence class of $r$ expressible through the above reparameterization. Define the projection $f$ as  
\begin{equation}
    f(r; \pi_\text{ref}, \beta)(x, y) = r(x, y) - \beta\log\sum_{y}\pi_\text{ref}(y\mid x)\exp\left(\frac{1}{\beta}r(x, y)\right)
\end{equation}
The operator $f$ acts by normalizing $r$ via subtracting the logarithm of the partition function (computed under $\pi_r$), which only depends on $x$. Consequently, $f(r; \pi_\text{ref}, \beta)(x, y)$ remains in the equivalence class of $r(x, y)$. Moreover, substituting $r$ using the expression from Eq.~\ref{eq:main_eq} (valid for any reward function), it follows that $f(r; \pi_\text{ref}, \beta)(x, y) = \beta \log \frac{\pi_r(y\mid x)}{\pi_\text{ref}(y\mid x)}$. Therefore, this projection yields a representative of the equivalence class for $r$ possessing the specified form, implying no loss of expressivity from adopting the proposed reparameterization of the reward model.
\end{sproof}

Theorem~\ref{thm:main} can also be interpreted as uniquely determining, within each equivalence class, the specific reward function that the DPO reparameterization chooses. This reward function must fulfill the following condition:

\begin{equation}\label{eq:lag_p}
     \sum_{y}\underbrace{\pi_\text{ref}(y\mid x)\exp\left(\frac{1}{\beta}r(x, y)\right)}_{=\pi(y\mid x)\text{, using Thm.~\ref{thm:main} reparam.}} = 1,
\end{equation}

meaning that $\pi(y\mid x)$ forms a proper probability distribution (with positive probabilities summing to 1). Furthermore, by referencing Eq.~\ref{eq:op_policy}, it becomes apparent that Eq.~\ref{eq:lag_p} corresponds to the partition function of the optimal policy that emerges from the reward function $r(x, y)$. The central idea underpinning the DPO algorithm is that it is possible to enforce specific constraints on the inherently under-specified Plackett-Luce family of preference models (including Bradley-Terry) in such a way that the expressive power of reward models is maintained. Simultaneously, these constraints allow for the analytic tractability of the optimal policy described in Eq.~\ref{eq:op_policy} for every input $x$.

\subsection{Instability of Actor-Critic Algorithms}
This framework can further be employed to analyze the instability observed in standard actor-critic methods applied in RLHF, such as PPO. Consistent with the RLHF process, we focus on the RL fine-tuning phase discussed in Section~\ref{section:prelims}. There are notable links between this setting and the control as inference paradigm \cite{levine2018reinforcement}, as formulated for the constrained RL problem in \ref{eq:RL}. Here, the agent’s policy $\pi_{\theta}(y\mid x)$ is parameterized, and training seeks to minimize the KL divergence $\mathbb{D}_{\text{KL}}[\pi_{\theta}(y|x) \mid \mid \pi^*(y\mid x)]$, where $\pi^*$, defined in Eq.~\ref{eq:optimum_model}, denotes the optimal policy derived from the reward function $r_{\phi}(y, x)$. After some manipulation, this yields the following optimization target:

\begin{equation}\label{eq:AC}
    \max_{\pi_{\theta}}\mathbb{E}_{\pi_{\theta}(y\mid x)}\bigg[\underbrace{r_{\phi}(x, y) -\beta\log\sum_{y}\pi_\text{ref}(y\mid x)\exp\left(\frac{1}{\beta}r_{\phi}(x, y)\right)}_{f(r_{\phi}, \pi_\text{ref}, \beta)} - \underbrace{\beta\log\frac{\pi_{\theta}(y\mid x)}{\pi_\text{ref}(y\mid x)}}_{\text{KL}}\bigg]
\end{equation}

The objective optimized here aligns with that of earlier studies 
\citep{ziegler2020finetuning, stiennon2022learning, bai2022training, ouyang2022training}, where a DPO-equivalent reward within the $r_{\phi}$ reward class was employed. Within this context, the normalization component in $f(r_{\phi}, \pi_\text{ref}, \beta)$ can be viewed as the soft value function associated with the reference policy $\pi_\text{ref}$. Although this term does not alter the ultimate optimizer's solution, omitting it would cause the policy gradient to exhibit high variance, potentially destabilizing the training process. One can incorporate the normalization factor by utilizing a learned value function; however, optimizing such a function may also pose practical challenges. Past work has often mitigated this issue by normalizing rewards relative to a human completion baseline, effectively using a single-sample Monte Carlo approximation for normalization. By contrast, the DPO reparameterization framework leads to a reward formulation that is inherently baseline-free. 

\section{Experiments}
This section presents a series of empirical investigations to assess DPO’s effectiveness at policy learning from preference data. Our initial experiments, carried out in a controlled text-generation scenario, seek to quantify how effectively DPO manages the trade-off between maximizing reward and minimizing KL-divergence with the reference policy, particularly in comparison to standard preference learning methods like PPO. Subsequently, we test DPO on larger-scale models and more challenging RLHF domains, such as summarization and conversational tasks. Our findings indicate that, even with minimal hyperparameter tuning, DPO achieves competitive or superior performance relative to strong baselines, such as RLHF implemented with PPO, and outperforms strategies that return the top trajectory out of $N$ samples using a learned reward. Prior to reporting these findings, we provide an overview of our experimental methodology; for further details, please refer to Appendix~\ref{app:exp_details}.

\textbf{Tasks.} Our study examines three distinct open-ended text generation tasks. In each case, models are trained to learn a policy based on a dataset of preferences denoted as $\mathcal{D}=\bigl\{x^{(i)}, y_w^{(i)}, y_l^{(i)}\bigr\}_{i=1}^N$. In the \textbf{controlled sentiment generation} task, $x$ represents the beginning of a movie review sourced from the IMDb dataset \cite{maas-EtAl:2011:ACL-HLT2011}, and the model is required to generate a continuation $y$ that expresses positive sentiment. To facilitate systematic evaluation, for this setting we automatically construct preference pairs for generations using an existing sentiment classifier, ensuring $p(\text{positive}\mid x,y_w)>p(\text{positive}\mid x,y_l)$. For SFT, GPT-2-large is fine-tuned until convergence using review samples from the IMDb train set (see App~\ref{app:sentiment_details} for more details). In the \textbf{summarization} task, $x$ refers to a Reddit forum post, and the goal is for the policy to produce a concise summary $y$ of the post's main content. We follow previous approaches and employ the Reddit TL;DR summarization dataset \citep{volske-etal-2017-tl}, supplemented by human preference data assembled by \citeauthor{stiennon2022learning}. Our SFT model is trained on summaries written by humans for forum posts,\footnote{\url{https://huggingface.co/CarperAI/openai_summarize_tldr_sft}} and RLHF optimization is conducted via the TRLX framework \citep{leandro_von_werra_2023_7790115}. The preference labels derive from human evaluations over outputs from a related, comparably fine-tuned SFT model, as curated by \citeauthor{stiennon2022learning}. For the \textbf{single-turn dialogue} task, $x$ consists of a user's prompt—ranging from scientific inquiries to personal advice requests—and the model must generate a reply $y$ that is both informative and appropriate. This uses the Anthropic Helpful and Harmless dialogue dataset \citep{bai2022training}, which comprises 170,000 dialogues between a human participant and an automated agent. Each dialogue concludes with a pair of candidate replies produced by a large (but unspecified) language model, annotated with human preferences for one response. Since there is no ready-made SFT model for this scenario, we construct one by fine-tuning a standard language model on the subset of completions marked as preferred.

\textbf{Evaluation.} We employ two distinct evaluation strategies in our experiments. To assess how well each algorithm optimizes the constrained reward maximization objective in a controlled sentiment generation scenario, we characterize each algorithm by plotting its achievable trade-off curve—specifically, the relationship between attained reward and KL-divergence from the reference policy. This analysis is feasible because the experiments utilize a ground-truth reward function, provided by a sentiment classifier. In contrast, practical applications rarely have access to the true reward function. Therefore, we gauge the performance of our algorithms in such real-world settings by measuring their \textit{win rate} compared to a baseline policy, leveraging GPT-4 as a stand-in for human judgments of summary quality and response helpfulness within the summarization and single-turn dialogue tasks, respectively. For the summarization task, baseline outputs are derived from the test set reference summaries, while for dialogue, the human-preferred response in the test data serves as the baseline. Although prior work has found that language models can outperform established metrics as evaluators \citep{Chen2023ExploringTU}, we supplement this with a human subject study to validate our reliance on GPT-4 for evaluation, as detailed in Sec.~\ref{sec:human-judgments}. Our results show that GPT-4's assessments exhibit high concordance with human ratings, with human-GPT-4 agreement matching or exceeding the agreement rate among independent human annotators.

\textbf{Methods.} Beyond DPO, our evaluations include multiple established techniques for tuning language models to follow human preferences. For simplicity, we experiment with zero-shot prompting using \textbf{GPT-J} \citep{gpt-j} in the summarization domain and 2-shot prompting with \textbf{Pythia-2.8B} \citep{biderman2023pythia} for dialogue. Further, we consider the \textbf{SFT} model and also \textbf{Preferred-FT}, which is fine-tuned via supervised learning on the selected completion $y_w$, sourced either from SFT (in controlled sentiment and summarization tasks) or from a generic LM (for single-turn dialogue). Another quasi-supervised method we evaluate is \textbf{Unlikelihood}~\citep{welleck2019neural}, which optimizes the model to increase the probability of $y_w$ while actively decreasing the probability assigned to $y_l$, with an optional `unlikelihood' weight parameter $\alpha\in[0,1]$. We also examine \textbf{PPO} \citep{schulman2017proximal}, utilizing a reward model trained on preference data, alongside \textbf{PPO-GT}, an oracle that uses the true reward function available only in the controlled sentiment context. Our sentiment-related experiments include two PPO-GT variants: a default implementation \cite{leandro_von_werra_2023_7790115}, and a custom version that incorporates reward normalization and additional hyperparameter optimization for improved results (these adjustments are likewise made when running PPO with learned rewards). Lastly, we evaluate the \textbf{Best of $N$} method, where $N$ outputs are sampled from the SFT model (or Preferred-FT in the case of dialogue), and the response with the highest score—according to a learned reward function from the preference data—is chosen. Although this approach demonstrates strong performance by separating reward model quality from PPO optimization, it becomes computationally expensive for even moderate values of $N$, as it requires $N$ samples per test query.

\subsection{How well can DPO optimize the RLHF objective?}

Standard RLHF algorithms employ a KL-constrained reward maximization objective, which serves to maximize the reward while simultaneously limiting the divergence of the policy from a reference policy. Consequently, an appropriate comparison of different methods must jointly consider both the achieved reward and the magnitude of the KL divergence; higher rewards accompanied by a substantial increase in KL divergence may be undesirable. Figure~\ref{fig:frontier-tldr-main} presents the trade-off curve between reward and KL for different approaches within the sentiment task. For each method, we perform several training trials, varying a policy conservativeness parameter in each run (target KL $\in\{3,6,9,12\}$ for PPO, $\beta \in \{0.05,0.1,1,5\}$, and $\alpha\in\{0.05,0.1,0.5,1\}$ for unlikelihood, with different random seeds for preferred-FT). In total, this results in 22 experiments. After every 100 training steps, up to convergence, we assess each policy using a set of test prompts, calculating both the mean reward based on the actual reward function and the average sequence-level KL\footnote{Defined as the total of the KL-divergences across all timesteps.} relative to the reference policy $\text{KL}\left(\pi\mid \mid \pi_\text{ref}\right)$. The results indicate that DPO achieves a markedly more favorable frontier, obtaining the highest rewards at comparatively low KL values. This outcome is significant for several reasons. Firstly, while both DPO and PPO optimize the same objective, DPO consistently exhibits greater efficiency, with its reward/KL profile surpassing that of PPO across the board. Secondly, DPO yields a superior frontier even compared to PPO when the latter is provided with access to ground truth rewards (PPO-GT).

\subsection{Can DPO scale to real preference datasets?}
\label{sec:dpo-real-datasets}
We proceed to assess the fine-tuning effectiveness of DPO on tasks involving summarization and single-turn dialogue. In the context of summarization, it is well-documented that commonly used automated metrics like ROUGE may show weak alignment with human judgments~\citep{stiennon2022learning}. Previous studies have demonstrated that optimizing LMs with PPO based on human feedback can result in more preferable summaries. To compare various methods, we generate outputs from models on the test partition of the TL;DR summarization dataset, then calculate the average win rate relative to ground-truth completions. Samples for all methods are drawn using a temperature sweep from 0.0 to 1.0, and the corresponding win rates are depicted in Figure~\ref{fig:frontier-tldr-main} (right). DPO, PPO, and Preferred-FT all use the identical GPT-J SFT model as a starting point\footnote{\url{https://huggingface.co/CarperAI/openai_summarize_tldr_sft}}. Our results indicate that DPO attains a win rate close to 61\% at temperature 0.0, surpassing PPO’s best performance of roughly 57\% at the same temperature. Furthermore, DPO achieves a higher maximum win rate compared to the best-of-$N$ reference baseline. It is important to highlight that the $\beta$ hyperparameter of DPO was not systematically tuned in these experiments, suggesting that actual performance may be underestimated. Additionally, DPO demonstrates significantly greater stability across varying sampling temperatures compared to PPO, whose effectiveness declines to match the baseline GPT-J model at higher temperatures. Preferred-FT offers negligible improvement over the base SFT model. Finally, a direct human evaluation of DPO and PPO, discussed in Section~\ref{sec:human-judgments}, shows that DPO outputs sampled at a temperature of 0.25 were favored over PPO outputs at the same temperature 58\% of the time.

For the single-turn dialogue setting, we assess various approaches using the portion of the Anthropic HH dataset test split \citep{bai2022training} that involves a single human-assistant exchange. When evaluating with GPT-4, the preferred completions from the test set serve as ground truth, allowing us to calculate win rates for each approach. Due to the absence of an established SFT benchmark for this domain, our methodology begins with a pre-trained Pythia-2.8B model. We then fine-tune a reference model using Preferred-FT on selected completions, ensuring the output distribution remains consistent, and proceed to train with DPO. Additional baselines include the strongest result from 128 Preferred-FT completions (noting that performance saturates at 128 for this particular evaluation; refer to Appendix Figure~\ref{fig:best-of-n}) as well as a 2-shot prompted Pythia-2.8B base model. Results show that DPO equals or surpasses competing methods at their respective optimal temperatures. We further analyze an RLHF model, trained with PPO on the Anthropic HH dataset \footnote{\url{https://huggingface.co/reciprocate/ppo_hh_pythia-6B}} and made available by a reputable source \footnote{\url{https://github.com/CarperAI/trlx/tree/main/examples/hh}}, but find that no prompt or sampling temperature yields better performance than the base Pythia-2.8B. Drawing on TL;DR findings and given both strategies optimize an identical reward signal, we interpret Best of 128 as an approximate indicator of PPO-level results. Taken together, DPO emerges as the only approach with practical computational efficiency that consistently improves upon the Anthropic HH dataset’s preferred completions, and it achieves results on par with or superior to the computationally intensive Best of 128 baseline. Lastly, Figure~\ref{fig:dialogue-main} illustrates that DPO attains optimal performance with relatively rapid convergence.

\subsection{Generalization to a new input distribution}

To more rigorously assess the robustness of PPO and DPO to distributional changes, we tested the policies learned during the Reddit TL;DR summarization task on an out-of-distribution set: specifically, the test split of news articles from the CNN/DailyMail corpus \citep{nallapati-etal-2016-abstractive}, utilizing the top-performing sampling temperatures from TL;DR (0 and 0.25). The outcomes of this evaluation are shown in Table~\ref{tab:ood}. We calculated GPT-4 win rates relative to the reference summaries in the datasets, employing the same GPT-4 (C) evaluation prompt as with Reddit TL;DR, but with the term ``forum post'' replaced by ``news article''. Within this novel distribution, DPO maintains a pronounced advantage over PPO. These findings offer preliminary support for the claim that DPO can match PPO in its ability to generalize to new domains, despite DPO’s lack of access to the supplementary unlabeled Reddit TL;DR prompts incorporated in PPO.

\subsection{Validating GPT-4 judgments with human judgments}
\label{sec:human-judgments}
To assess the credibility of GPT-4 as an evaluator, we carried out a human annotation study, focusing on the TL;DR summarization results with two variations of GPT-4 prompts. The \textbf{GPT-4 (S)} (simple) prompt directly queries which summary better captures the essential content of the post, while the \textbf{GPT-4 (C)} (concise) prompt includes an explicit emphasis on conciseness; we analyze this prompt variant as we observe that with the \textbf{GPT-4 (S)} prompt, GPT-4 tends to prefer longer and more redundant summaries compared to human raters. Complete prompt text can be found in Appendix~\ref{app:prompts}. Three sets of comparisons are conducted: the highest-performing method (DPO, temp. 0.25), the lowest (PPO, temp. 1.0), and an intermediate model (SFT, temp. 0.25), each pitted against PPO with greedy sampling (its optimal temperature), in order to span a spectrum of summary quality. In both prompting conditions, agreement between GPT-4 and human judgments is similar to the agreement among different human annotators themselves, indicating that GPT-4 serves as a suitable surrogate for human preference assessments (due to limited human participation, we obtained multiple human judgments only for the DPO and PPO-1 pairwise comparisons). Overall, win rates using the \textbf{GPT-4 (C)} prompt most closely match those of human raters; for this reason, Section~\ref{sec:dpo-real-datasets} relies primarily on results from the \textbf{GPT-4 (C)} prompt. Further methodological specifics, including details on the annotation platform and the list of participating human raters, are presented in Appendix~\ref{app:human-study}.

\section{Discussion}
Preference-based learning offers a robust and scalable method for developing language models that are both capable and aligned. In this work, we present DPO, a straightforward approach to training language models from preference data that does not rely on reinforcement learning. Unlike prior methods that adapt the preference learning framework to fit conventional RL formulations and thereby employ standard RL algorithms, DPO directly links language model policies to reward functions. This allows for the direct training of a language model to align with human preferences using a standard cross-entropy loss, eliminating the need for reinforcement learning and maintaining general applicability. DPO achieves competitive or superior results compared to established RLHF methods such as PPO, even with minimal hyperparameter adjustment, effectively lowering the technical barriers associated with preference-based model training.

\textbf{Limitations \& Future Work.} This research opens up multiple avenues for further investigation. One question pertains to how well the DPO policy transfers to out-of-distribution cases compared to approaches that rely on explicit reward modeling; while preliminary findings indicate that DPO's generalization abilities are on par with PPO-based alternatives, further systematic evaluation is necessary. Another direction is to assess whether self-labeling using the DPO policy can capitalize on unlabeled prompts as effectively as supervised approaches. Additionally, the way reward over-optimization may present itself in the DPO framework is an open question—whether the marginal dip in results shown in Figure~\ref{fig:dialogue-main}-right reflects this phenomenon requires deeper inquiry. Though our evaluation encompasses models of up to 6B parameters, scaling DPO for models substantially larger than this remains an open challenge for future exploration. With regard to assessment, our findings highlight that win rates obtained via GPT-4 are sensitive to prompt variations; subsequent work could aim to determine more reliable methods for gathering high-quality automated evaluations. Beyond training language models to adhere to human preferences, DPO holds potential for a range of applications, such as in the training of generative models across other domains.